\begin{lemma}[Termination and a tight residual cut]
Every finite integer-capacity network with distinct terminals admits a
feasible integral flow $f$ and a subset $S$ containing $s$ and excluding
$t$ such that
\[
 |f|=C(S),\qquad 0\leq |f|\leq\sum_w c(s,w).
\]
Here network, flow, and cut capacity have the meanings in
\noderef{FiniteCapacityNetwork} and \noderef{FiniteIntegralFlow}.
\end{lemma}
\begin{proof}
The capacities in \noderef{FiniteCapacityNetwork} are finite nonnegative
integers. The zero function satisfies \noderef{FiniteIntegralFlow}: it is
skew symmetric, is bounded above by every capacity, and has zero divergence.
For any feasible flow, summing its upper bounds at the source gives
$|f|\leq M:=\sum_w c(s,w)$, a finite nonnegative integer.

Start at zero. If a simple residual source-to-sink path exists, apply
\noderef{IntegralFlowAugmentation}. This preserves feasibility and
increases the integer value by at least one. This process must stop
within $M$ augmentations: after $r$ augmentations its value is at least
$r$, and a further augmentation from value $M$ contradicts the bound.
More formally, induction on the nonnegative integer $M-|f|$ gives a
terminal flow reachable from any intermediate flow: without a path it
is already terminal, and with a path augmentation decreases this integer
strictly. Starting at zero therefore gives a terminal $f$ with
$0\leq |f|\leq M$ and no simple residual source-to-sink path.

Let $S$ consist of vertices reachable from $s$ by finite residual walks,
including the length-zero walk. It is a finite subset because the vertex
type is finite, and it contains $s$. A residual walk from $s$ to $t$
would give a simple path: if a vertex repeats at positions $i<j$, delete
the segment after position $i$ through position $j$. The adjoining step
is still a residual step because the vertices at $i,j$ agree. This
shortens the walk and preserves its endpoints. Repetition removal
terminates as lengths decrease, producing an injective walk. Distinct
terminals ensure its length is positive. It is therefore a path in the
precise sense of \noderef{SimpleResidualPath}, contrary to terminality.
Thus $t\notin S$.

For $v\in S,w\notin S$ there is no residual step from $v$ to $w$,
since appending such a step would make $w$ reachable. Consequently
$f(v,w)\geq c(v,w)$, and the flow bound gives equality. Apply
\noderef{IntegralFlowCutIdentity} and replace each crossing net flow
by its capacity. This gives $|f|=C(S)$ as required. When $M=0$ the same
argument starts and stops at the zero flow; no positive-capacity path
can be augmented without contradicting the value bound.
\end{proof}
